% GENERATED FILE. Edit canonical lessons, then run npm run generate:books.
% canonical-source-hash: fnv1a64:d2d3f6d69212a289
% canonical-lessons: RU-C237-umirat, RU-C237-stuchat, RU-C237-remontirovat, RU-C237-pereyezzhat, RU-C237-zarabatyvat

\chapter{To die, To knock, To renovate, To move, To earn}
\label{ch:ru-to-die-to-knock-to-renovate-to-move-to-earn}
\hlchaptermodality{237}

\begin{chapteropening}
\textbf{By the end of this chapter:} \emph{I can say to die, to knock, to renovate, to move and to earn.}

Complete the last lesson of chapter 237: I can say to die, to knock, to renovate, to move and to earn.
\end{chapteropening}

\section[\texorpdfstring{\ru{умирать} (umirát)}{умирать (umirát)}]{\ru{умирать} (umirát) — to die}

\label{lesson:RU-C237-umirat}



\begin{quote}
Before the new one: say the Russian for to check, then the Russian for to fill in.
\end{quote}

\subsection*{You'll want to know: \ru{умирать}}
\textbf{\ru{умирать}} (\emph{umirát}) — \textquotedblleft{}to die\textquotedblright{}.

\begin{cousinweb}
\textbf{\ru{умирать}} is \emph{\ru{у}-} with \emph{-\ru{мирать}}, from \emph{\ru{мереть}}, \textquotedblleft{}to die\textquotedblright{}, Proto-Slavic \emph{*merti}, which Vasmer relates to Lithuanian \emph{mìrti} and Latin \emph{morī}, \textquotedblleft{}to die\textquotedblright{}.
\end{cousinweb}

\begin{grammarlens}[title={What you've built}]
One more word for this chapter.
\end{grammarlens}

\subsection*{Your turn}
\begin{itemize}
\raggedright
  \item \emph{Say it:} \emph{umirát}
  \item \emph{Say it:} \emph{umirát}, once more
  \item \emph{Say it:} say \emph{umirát}
  \item \emph{Recall:} say the Russian for to check, then the Russian for to fill in, then say \emph{umirát} again
\end{itemize}

\begin{tcolorbox}[breakable,colback=teal!4,colframe=teal!35!black,title={Before you move on}]
What does \ru{умирать} mean? (\textquotedblleft{}To die\textquotedblright{}.) Say it once more.
\end{tcolorbox}
\section[\texorpdfstring{\ru{стучать} (stuchát)}{стучать (stuchát)}]{\ru{стучать} (stuchát) — to knock}

\label{lesson:RU-C237-stuchat}



\begin{quote}
Before the new one: say the Russian for to fill in, then the Russian for to die.
\end{quote}

\subsection*{You'll want to know: \ru{стучать}}
\textbf{\ru{стучать}} (\emph{stuchát}) — \textquotedblleft{}to knock\textquotedblright{}.

\begin{cousinweb}
\textbf{\ru{стучать}} comes from \emph{\ru{стук}}, \textquotedblleft{}a knock\textquotedblright{} (Vasmer), which he compares with sound words like Czech \emph{ťukati}, \textquotedblleft{}to tap\textquotedblright{}, and French \emph{toc}.
\end{cousinweb}

\begin{grammarlens}[title={What you've built}]
One more word for this chapter.
\end{grammarlens}

\subsection*{Your turn}
\begin{itemize}
\raggedright
  \item \emph{Say it:} \emph{stuchát}
  \item \emph{Say it:} \emph{stuchát}, once more
  \item \emph{Say it:} say \emph{stuchát}
  \item \emph{Recall:} say the Russian for to fill in, then the Russian for to die, then say \emph{stuchát} again
\end{itemize}

\begin{tcolorbox}[breakable,colback=teal!4,colframe=teal!35!black,title={Before you move on}]
What does \ru{стучать} mean? (\textquotedblleft{}To knock\textquotedblright{}.) Say it once more.
\end{tcolorbox}
\section[\texorpdfstring{\ru{ремонтировать} (remontírovat)}{ремонтировать (remontírovat)}]{\ru{ремонтировать} (remontírovat) — to renovate, to repair}

\label{lesson:RU-C237-remontirovat}



\begin{quote}
Before the new one: say the Russian for to die, then the Russian for to knock.
\end{quote}

\subsection*{You'll want to know: \ru{ремонтировать}}
\textbf{\ru{ремонтировать}} (\emph{remontírovat}) — \textquotedblleft{}to renovate, to repair\textquotedblright{}.

\begin{cousinweb}
\textbf{\ru{ремонтировать}} is \emph{\ru{ремонт}}, \textquotedblleft{}repair\textquotedblright{}, made into a verb. \emph{\ru{ремонт}} comes from French \emph{remonte}, \textquotedblleft{}a fresh supply of horses\textquotedblright{} (Vasmer); under Peter the Great it meant supplying horses to regiments.
\end{cousinweb}

\begin{grammarlens}[title={What you've built}]
One more word for this chapter.
\end{grammarlens}

\subsection*{Your turn}
\begin{itemize}
\raggedright
  \item \emph{Say it:} \emph{remontírovat}
  \item \emph{Say it:} \emph{remontírovat}, once more
  \item \emph{Say it:} say \emph{remontírovat}
  \item \emph{Recall:} say the Russian for to die, then the Russian for to knock, then say \emph{remontírovat} again
\end{itemize}

\begin{tcolorbox}[breakable,colback=teal!4,colframe=teal!35!black,title={Before you move on}]
What does \ru{ремонтировать} mean? (\textquotedblleft{}To renovate, to repair\textquotedblright{}.) Say it once more.
\end{tcolorbox}
\section[\texorpdfstring{\ru{переезжать} (pereyezzhát)}{переезжать (pereyezzhát)}]{\ru{переезжать} (pereyezzhát) — to move (house)}

\label{lesson:RU-C237-pereyezzhat}



\begin{quote}
Before the new one: say the Russian for to knock, then the Russian for to renovate.
\end{quote}

\subsection*{You'll want to know: \ru{переезжать}}
\textbf{\ru{переезжать}} (\emph{pereyezzhát}) — \textquotedblleft{}to move (house)\textquotedblright{}.

\begin{cousinweb}
\textbf{\ru{переезжать}} is \emph{\ru{пере}-}, \textquotedblleft{}across\textquotedblright{}, with \emph{-\ru{езжать}}, from \emph{\ru{ездить}}, \textquotedblleft{}to travel\textquotedblright{}: to move from one home to another.
\end{cousinweb}

\begin{grammarlens}[title={What you've built}]
One more word for this chapter.
\end{grammarlens}

\subsection*{Your turn}
\begin{itemize}
\raggedright
  \item \emph{Say it:} \emph{pereyezzhát}
  \item \emph{Say it:} \emph{pereyezzhát}, once more
  \item \emph{Say it:} say \emph{pereyezzhát}
  \item \emph{Recall:} say the Russian for to knock, then the Russian for to renovate, then say \emph{pereyezzhát} again
\end{itemize}

\begin{tcolorbox}[breakable,colback=teal!4,colframe=teal!35!black,title={Before you move on}]
What does \ru{переезжать} mean? (\textquotedblleft{}To move (house)\textquotedblright{}.) Say it once more.
\end{tcolorbox}
\section[\texorpdfstring{\ru{зарабатывать} (zarabátyvat)}{зарабатывать (zarabátyvat)}]{\ru{зарабатывать} (zarabátyvat) — to earn}

\label{lesson:RU-C237-zarabatyvat}



\begin{quote}
Before the new one: say the Russian for to renovate, then the Russian for to move.
\end{quote}

\subsection*{You'll want to know: \ru{зарабатывать}}
\textbf{\ru{зарабатывать}} (\emph{zarabátyvat}) — \textquotedblleft{}to earn\textquotedblright{}.

\begin{cousinweb}
\textbf{\ru{зарабатывать}} is \emph{\ru{за}-} with \emph{-\ru{рабатывать}}, from \emph{\ru{работать}}, \textquotedblleft{}to work\textquotedblright{}: to earn is to get by working.
\end{cousinweb}

\begin{grammarlens}[title={What you've built}]
One more word for this chapter.
\end{grammarlens}

\subsection*{Your turn}
\begin{itemize}
\raggedright
  \item \emph{Say it:} \emph{zarabátyvat}
  \item \emph{Say it:} \emph{zarabátyvat}, once more
  \item \emph{Say it:} say \emph{zarabátyvat}
  \item \emph{Recall:} say the Russian for to renovate, then the Russian for to move, then say \emph{zarabátyvat} again
\end{itemize}

\begin{tcolorbox}[breakable,colback=teal!4,colframe=teal!35!black,title={Before you move on}]
What does \ru{зарабатывать} mean? (\textquotedblleft{}To earn\textquotedblright{}.) Say it once more.
\end{tcolorbox}
